%==============================
\subsection{Overall Analyses}\label{subsec:overall}

\subsubsection{Execution and Online Adaptation}\label{subsec:execute}
The hybrid search returns a pushing schedule
$\pi = \{\tau_1,\cdots,\tau_K\}$, where
$\tau_k \triangleq (\Omega_k,\mathbf{v}_k,\boldsymbol{\xi}_k)$ specifies
the target obstacle, pushing direction and mode. The schedule is executed
sequentially in a receding-horizon manner. For each $\tau_k$, the robots move
to the required mode~$\boldsymbol{\xi}_k$ for obstacle~$\Omega_k$
and follow the pushing direction~$\mathbf{v}_k$
from the current state as in~\eqref{eq:sim}. On the hardware, optimized contact forces are used for mode feasibility rather than direct force commands; the controller converts the planned contact geometry and object motion into velocity commands with feedback tracking. During transition, robots
are reassigned by preserving the circular order of contact points along the
obstacle boundary~\cite{tang2026pushingbots}, which reduces path crossings
under sufficient clearance. After each push, the system state is updated from
observation, and the execution proceeds to $\tau_{k+1}$ if the action succeeds.
For consistency with execution, the WCCG is rebuilt periodically from
the current configuration. Execution terminates once a $W$-clear path
$\mathcal{P}^W_\texttt{V}$ connects~$\mathbf{x}_\texttt{V}^{\texttt{S}}$ and
$\mathbf{x}_\texttt{V}^{\texttt{G}}$.
Replanning is triggered by observed execution discrepancies, including failed contact establishment, stalled pushing, or invalidated remaining WCCG connectivity. In such cases, Alg.~\ref{alg:pihs} is
called from the current state to derive the updated pushing schedule.

%----------------------------------------------
\subsubsection{Computational Complexity}\label{subsubsec:complexity}

Given~$M$ movable obstacles, the computational cost of one node
expansion in Alg.~\ref{alg:pihs} is dominated by candidate generation and
physics validation. The WCCG update and gap ranking incur
$\mathcal{O}(M\log M+|\mathcal{V}|+|\mathcal{E}_c|+|\mathcal{E}_b|)$ and
$\mathcal{O}(|\Gamma_{\mathcal{L}_\nu}|
\log|\Gamma_{\mathcal{L}_\nu}|)$ costs, respectively, while the dominant
simulation cost is
\[
\mathcal{O}\left(
\sum_{g\in\mathsf{Rank}(\nu)}
|\Xi_g|
\left(
C_{\mathrm{LP}}+
\frac{\rho_\nu T_{\mathrm{sim}}}{n_{\texttt{w}}}
\right)
\right),
\]
Therefore, scalability is mainly determined by validated pushing candidates,
while quick-pass filtering and parallel simulation reduce the practical
expansion cost. Robot-team size mainly affects mode generation, where greedy
contact refinement avoids exhaustive $|\mathcal{A}|^N$ contact assignment and
evaluates bounded single-contact replacements.
%% %----------------------------------------------
%% \subsubsection{Generalization}\label{subsec:general}
%% The framework admits several extensions:
%% (I) \textit{Heterogeneous teams.} Per-robot gain tuning and robot-specific mode
%% priors enable cooperation among diverse robots;
%% (II) \textit{Concurrent pushing.} Multi-object mode generation and conflict-aware
%% transition planning allow multiple obstacles to be pushed simultaneously;
%% (III) \textit{Dynamic environments.} Periodic $W$--connectivity checks and
%% on-demand replanning handle disturbances such as object drift or unmodeled
%% contacts during execution.
